\def\ProblemCode{DSUROLLBACK}
\def\ProblemTitle{DSU Rollback (Basic)}
\makeproblem{\ProblemCode}{\ProblemTitle}

Cho đồ thị vô hướng gồm $n$ đỉnh. Các đỉnh được đánh số từ $1$ tới $n$. Ban đầu, đồ thị chưa có cạnh nào và các đỉnh $1,2,\dots,n$ lần lượt mang giá trị $a_1,a_2,\dots,a_n$.

Hãy xử lý $q$ truy vấn thuộc một trong các dạng sau:

\begin{itemize}
    \item \boxed{0\ u\ v}: Thêm một cạnh mới nối giữa hai đỉnh $u$ và $v$.
    \item \boxed{1}: Xóa cạnh gần nhất đã được thêm vào từ một thao tác loại \texttt{0} mà chưa bị xóa.
    \item \boxed{2}: Đếm số thành phần liên thông trong đồ thị.
    \item \boxed{3\ u}: Đếm số đỉnh thuộc cùng thành phần liên thông với đỉnh $u$ (bao gồm cả $u$).
    \item \boxed{4\ u}: Tìm giá trị đỉnh nhỏ nhất trong các đỉnh thuộc cùng thành phần liên thông với $u$.
\end{itemize}

\textbf{Lưu ý:} Giữa hai đỉnh có thể có nhiều cạnh nối phân biệt.

% ===== INPUT DESCRIPTION =====
\inputDesc{}

\begin{itemize}
    \item Dòng đầu tiên gồm hai số nguyên dương $n, q\ (1 \le n,q \le 10^{5})$.
    \item Dòng tiếp theo gồm $n$ số nguyên $a_1,a_2,\dots,a_n\ (|a_i| \le 10^{9})$.
    \item $q$ dòng tiếp theo, mỗi dòng chứa một truy vấn.
\end{itemize}

Dữ liệu đảm bảo tại mọi thời điểm, số truy vấn loại \texttt{0} xuất hiện không ít hơn số truy vấn loại \texttt{1}.

% ===== OUTPUT DESCRIPTION =====
\outputDesc{}

\begin{itemize}
    \item Với mỗi truy vấn loại 2, 3 hoặc 4, in ra một dòng chứa kết quả tương ứng.
\end{itemize}

% ===== SUBTASKS =====
% \noindent{\sffamily \textbf{Subtasks}}
% \begin{itemize}
%     \item \textbf{Subtask ?} ($?\%$ số điểm): Không có ràng buộc gì thêm.
% \end{itemize}

\vspace{0.5em}

\ExampleBox{1}{
5 8\\
3 -1 4 2 0\\
2\\
0 1 2\\
4 1\\
0 2 3\\
3 1\\ 
1\\
2\\
4 3
}{
5\\
-1\\
3\\
4\\
4
}{}
